\documentclass[12pt,a4paper]{article}
\usepackage[UTF8]{ctex}
\usepackage{amsmath,amssymb,graphicx,booktabs,caption,geometry,xcolor}
\geometry{left=2.5cm,right=2.5cm,top=2.5cm,bottom=2.5cm}

\begin{document}

\section{问题三的模型建立与求解：基于信息驱动的动态搜索与连锁清除策略}

\subsection{问题分析与算法引入}
问题三的核心挑战在于：在未知干扰源数量（$10\sim 16$ 个）和位置的黑盒环境下，控制机器狗在半径 1800 米的大圆区域内完成搜索与清除，且总耗时极小化。系统的总耗时由移动时间（5米/秒）、频道切换时间（1秒/次）、检测时间（5秒/次）、光学精确定位时间（3秒）及清除时间（2秒）构成。

如果采用传统的“全图网格化固定路线扫描”并在最后统一清除，将不可避免地产生大量重复的移动路径和冗余的空扫时间。为此，本文将该问题建模为带有不完全信息的动态车辆路径规划问题（Dynamic VRP with Partial Observability），并引入\textbf{贪心算法（Greedy Algorithm）}的思想来主导机器狗的在线决策。
\textbf{贪心算法简介}：贪心算法是指在对问题求解时，总是做出在当前看来是最好的选择。在本文的语境下，这意味着机器狗在每一个决策节点，都会优先选择距离最近、且能够立刻被清除的干扰源作为下一个目标。尽管单步贪心算法存在陷入局部最优的风险（例如为了清除近处的一个孤立目标而偏离了目标密集的集群区域），但通过引入多步前瞻（N-Step Look-ahead）或结合空间状态剪枝，我们能将其打造成一个极具实效性的动态寻优引擎。

为了将总耗时压缩至极限，本文在考虑 $\pm 1^\circ$ 测向误差与定位精度的前提下，构建了包含“状态机管理”、“动态频道剪枝”、“连锁探测”以及“偏航逼近”的综合算法模型。

\subsection{严格误差约束下的干扰源状态机模型}
定义全体 $N=20$ 个可选频道的集合，为每个频道赋予五种严格受控的动态状态 $S_i$：
\begin{itemize}
    \item \textbf{状态 0（未知，Unseen）}：尚未被探测到。
    \item \textbf{状态 1（单线弱定位，1-Ray）}：仅获得1条射线，无法形成交会。
    \item \textbf{状态 2（粗定位，Rough）}：获得 $\ge 2$ 条射线，但根据 $\pm 1^\circ$ 的测向误差计算，其交叉区域的最小包围圆半径 $\rho > 20$ 米。此时机器狗不能直接前往目标点，否则无法触发光学定位。
    \item \textbf{状态 3（精确定位，Precise）}：交会误差 $\rho \le 20$ 米，目标已彻底暴露，具备光学清除条件。
    \item \textbf{状态 4（已清除，Eliminated）}：完成清除操作，永久移出任务列表。
\end{itemize}

其中交叉包围圆半径 $\rho$ 的近似估算公式为：
\begin{equation}
    \rho \approx \frac{\max(d_1, d_2) \cdot (\pi/180)}{\sin \alpha}
\end{equation}
公式表明：当探测距离 $d$ 过大或交汇角 $\alpha$ 过小时，定位误差将急剧放大。

\subsection{空间频段动态剪枝策略（时间极限压缩法一）}
在任意位置进行全频段扫描时，盲目扫描所有未清除的频道将耗费大量时间（每多扫一个频道耗时 6 秒）。本文提出\textbf{基于空间位置的频道动态过滤算法}：
在机器狗进行探测前，模型会计算当前位置与所有已获知粗略位置（状态1或2）干扰源的理论距离。若当前位置距离某干扰源的估算区域下界 $>1000$ 米（超出了题设的有效接收半径下限），则可以断定在当前位置绝对无法接收到该频道的信号。
此时，算法在本次检测的频道轮询列表中\textbf{临时剔除（剪枝）}该频道。此举完美避免了无效的等待时间，极大地缩减了每次驻留时的无意义检测耗时。

\subsection{机器狗动态路由控制与行动决策机制}
算法初始时，规划 7 个全局粗探点（原点及外围半径 1150 米的正六边形顶点，确保 1800 米区域全局覆盖）。机器狗的行为逻辑由状态机触发，按以下优先级执行动作：

\noindent\textbf{优先级 1：局部清剿与“连锁清除反应”}。若存在状态 3 的频道，机器狗运用贪心算法，立即飞向距离最近的该目标点完成清除（转为状态 4）。
\textbf{核心优化（时间极限压缩法二）}：完成清除后，机器狗\textbf{不立即离开}，而是在该清除点原位作为“新观测基站”进行下一次探测。由于干扰源分布往往存在局部聚集性，这一操作不仅获得了极佳的近距离信号，更能与其他基站的射线形成交叉，顺带确立附近其他干扰源的位置，从而引发“扫雷”般的高效连锁清除反应。

\noindent\textbf{优先级 2：打破“射线共线陷阱”的横向偏航逼近机制}。若无状态 3 目标，但存在状态 2（粗定位 $\rho > 20$ 米）目标。如果机器狗沿着原有射线直直飞向目标，新旧射线将发生重合（共线，即 $\alpha \approx 0$），导致式(1)中的误差无限放大，使算法陷入死循环。
\textbf{核心优化（破局策略）}：引入偏航逼近机制。机器狗不盲目直飞，而是引入 $30^\circ$ 的偏航角进行斜向逼近，并在距离目标估算点 250 米处作为“间接中间探测点”停下。在 250 米处，$\pm 1^\circ$ 的横向误差仅约 4.3 米，下一次扫描必定能以较大的夹角产生 $\rho \le 20$ 的完美交叉。若距离极近但仍处于状态 2，则强制机器狗进行横向切向移动 50 米，人为拉大交叉基线。

\noindent\textbf{优先级 3：全局探索开图}。若已知信号均不足以引导逼近，机器狗按旅行商问题（TSP）求得的最短顺序前往下一个全局探测点获取全新射线，开辟新视野。

\subsection{模型异常处理与备选极值策略（附加扩展）}
为了提高模型在极端恶劣电磁环境下的鲁棒性，本文还提出了以下备选极值策略：
\begin{enumerate}
    \item \textbf{极小区域网格盲搜}：当交会误差被压缩至 $30\sim 40$ 米，但受限于地形无法进一步拉开基线测向时，直接舍弃数学交会。利用光学探测仪 20 米的有效范围，在 40 米的盲区内执行“Z字型”物理巡场，只需移动极短距离即可触发光学精确定位。
    \item \textbf{基于场强衰减的单站在线自适应标定}：利用题目“场强仅随距离衰减”的特性，采用对数距离路径损耗模型 $RSSI(d) = A - 10 \gamma \log_{10}(d)$。在清除首个目标的过程中收集远近两次场强数据，在线反解标定出环境参数 $A$ 和 $\gamma$。此后，对于新目标，只需一次探测即可通过场强反推距离 $d$，结合示向角度实现单站极坐标确切定位，彻底打破多点交会的时间瓶颈。
\end{enumerate}

\end{document}
